% Append-only section for the AION technical record.
% Requires: amsmath, booktabs, url.
\section{Transfer to Qwen3-30B-A3B on an 18~GB Commodity Mac}
\label{sec:qwen30b-transfer-2026-09-09}

\subsection{Objective and claim boundary}

This stage tests whether the storage-addressing and bounded-execution principles
developed on Granite transfer to a substantially larger sparse mixture-of-experts
(MoE) model.  The target is the official Qwen3-30B-A3B Q4\_K\_M GGUF checkpoint,
stored on the removable SD evidence volume and executed on the existing Apple M3
Pro Mac with 18~GB unified memory.  The aspirational performance target is
50 generated tokens per second, but it is neither assumed nor used to move an
acceptance threshold.  At the end of this stage, feasibility and the storage
addressing mechanism are established; a high-speed Qwen result is not yet
claimed.

\subsection{Storage preparation and provenance}

Six historical Granite layer-pack directories, occupying approximately 31~GiB,
were moved intact from the SD card to the reversible internal archive
\texttt{/Users/kevinrobinson/AION-Artifact-Archive/2026-09-09-granite-layer-packs}.
The relocation manifest records the six source and destination hashes, all of
which matched after relocation.  The original Granite model, promoted expert
banks, expert blocks, experiment evidence, runtime evidence and technical records
were retained on the SD volume.  This increased available SD capacity from
approximately 5.4~GiB to 36~GiB before the Qwen checkpoint was installed.

The installed checkpoint has the following immutable identity:
\begin{align*}
  B_{\mathrm{model}} &= 18{,}556{,}685{,}824\ \mathrm{bytes},\\
  H_{\mathrm{model}} &= \texttt{0d003f6662faee786ed5da3e31b29c97}\ldots
                       \texttt{c606a7f3c01aa8f5}.
\end{align*}
The digest matched the upstream linked-file identity.  A complete sequential
SHA-256 pass required 194.14~s, corresponding to 95.58~MB/s decimal sustained
read throughput from the SD card.  This provides a directly measured cold
storage floor rather than an advertised device rate.  The runtime was
\texttt{llama.cpp 0.4.0}, build 10809, commit \texttt{5266f24da}.

\subsection{Failed stock Metal control}

The first completed stock placement used all eligible GPU layers, kept MoE
weights on CPU, enabled memory mapping and retained the runtime's default weight
repacking.  It generated no token and ended after 574.07~s with repeated Metal
\texttt{Insufficient Memory} command-buffer failures.  This run is preserved
with an \texttt{INVALID} sidecar and supplies no tokens-per-second value.  It
demonstrates that storing the checkpoint is possible but conventional
whole-checkpoint materialisation is not a viable execution path on this machine.

Several measurement-harness defects were also caught and preserved: one rejected
argument, one pseudo-terminal job-control stop, one chat frontend that wrote
approximately 90~MB of spinner characters to the SD result stream, and one
unsupported benchmark argument.  None is used as performance evidence.  The
accepted control uses the non-interactive batched benchmark and compact JSONL
output.

\subsection{Valid CPU/no-repack feasibility baseline}

The accepted control disabled repacking, assigned zero layers to Metal, retained
all MoE tensors on CPU/SD, used memory mapping, and fixed context 1024,
batch/physical batch 64, eight CPU threads, 32 prompt tokens, eight generated
tokens and one sequence.  It completed without an out-of-memory failure.

\begin{table}[ht]
\centering
\begin{tabular}{lr}
\toprule
Metric & Measured value \\
\midrule
Cold wall time & 358.55~s \\
Prompt time (32 tokens) & 44.878975~s \\
Prompt throughput & 0.713029 tokens/s \\
Generation time (8 tokens) & 14.593630~s \\
Generation throughput & 0.548184 tokens/s \\
Combined throughput (40 tokens) & 0.672579 tokens/s \\
\bottomrule
\end{tabular}
\caption{First valid Qwen3-30B-A3B execution baseline on the SD-backed 18~GB M3 Pro.}
\end{table}

The result artifact SHA-256 is
\texttt{e2bdc67c7df654c1d079f4f997f95a14e4bae8de11a49287d7c425c26b002e06};
the runtime-log SHA-256 is
\texttt{bfe862576d22a0a1578d6a6eff90140bf6257b86eb37174fa5fdf70543fe60af}.
The benchmark establishes computational feasibility, not answer quality, because
the batched harness measures token identifiers rather than a semantic response.

\subsection{Whole-GGUF Metal placement is falsified}

Disabling repacking did not make conventional Metal placement safe.  Placing all
eligible layers on Metal crossed Apple's reported 12{,}884.92~MB recommended
working set and terminated abnormally after 287.46~s.  A 24-layer bound crossed
the same limit and terminated after 286.29~s.  An eight-layer bound also crossed
the limit, reported Metal out-of-memory during decode at approximately 8~min
29~s, and failed the fixed 600-second wall gate.  None emitted valid benchmark
timing.

This falsifies the idea that a small adjustment to the number of offloaded layers
is sufficient.  On Apple silicon, CPU mappings and Metal allocations compete for
one unified-memory pool.  The complete GGUF mapping must therefore be prevented
from becoming the effective Metal residency unit.  Future execution must expose
only bounded, currently required layer/expert buffers.

\subsection{Zero-copy Glyph expert address map}

GGUF inspection found 579 tensors.  Each of the 48 MoE layers stores its 128
experts in three combined tensors: gate, up and down projections.  Each combined
tensor divides evenly into 128 physically contiguous expert slices.  AION now
parses the GGUF header and constructs a source-bound address for every slice
without reading or copying the tensor payload.

For layer $\ell$, expert $e$ and projection $p$, the address is
\begin{equation}
  G(\ell,e,p) = (H_{\mathrm{model}}, o_{\ell,p} + e s_{\ell,p},
                 s_{\ell,p}, \tau_{\ell,p}),
\end{equation}
where $o$ is the absolute tensor offset, $s$ is its per-expert byte stride and
$\tau$ records the GGML type and tensor dimensions.  The human-readable Glyph is
\texttt{qwen3moe/layer/LL/expert/EEE}.

The generated map contains 6{,}144 layer/expert instances and has these measured
properties:
\begin{align*}
  B_{\mathrm{expert}} &= 17{,}553{,}162{,}240\ \mathrm{bytes},\\
  \frac{B_{\mathrm{expert}}}{B_{\mathrm{model}}} &= 94.5921\%,\\
  B_{\mathrm{one\ layer/expert}} &= 2{,}856{,}960\ \mathrm{bytes},\\
  B_{\mathrm{eight\ active,48\ layers}} &= 1{,}097{,}072{,}640\ \mathrm{bytes/token}.
\end{align*}
The last quantity is the logical uncached expert demand, not necessarily physical
SD traffic after reuse.  It proves why naive streaming cannot approach the target:
at the measured 95.58~MB/s SD rate, reading 1.097~GB for every token would itself
take more than eleven seconds before computation.  Large gains require reuse,
bounded residency and direct packed computation, not merely more reader threads.

The 6.6~MB address-map file SHA-256 is
\texttt{70e3ca7f0db99a9f357a5d3789a56b4df878b667b1d0b2978439f4b317bf6ddc};
its internal canonical digest is
\texttt{f9518b63658231ccad3d07d48a445aa7303cf1fc29148bf4bb13f8aabda4e82e}.
The parser and synthetic coverage test were committed in isolated commit
\texttt{0ac48511}; focused result: 1/1 passed.

\subsection{Bounded real-SD expert-store mechanism gate}

A strict 256~MiB least-recently-used store was then connected to the address map.
It reads the three exact packed ranges for a requested layer/expert pair using
positional reads, rejects short reads, retains no more than its declared byte
ceiling and never copies or rewrites the source model.  The first report was
invalidated because it estimated logical demand using a model-wide mean expert
size; quantized expert sizes vary slightly by layer.  The corrected report sums
the exact declared byte ranges for every access.

The narrow mechanism trace comprised four synthetic token routes over layers
0--7, with eight unique experts per layer and controlled partial overlap.  It is
not a model-observed Qwen route, so it cannot establish production cache hit rate
or tokens per second.  The corrected measured result was:
\begin{align*}
 B_{\mathrm{logical,no\ cache}} &= 757{,}334{,}016\ \mathrm{bytes},\\
 B_{\mathrm{physical}} &= 568{,}000{,}512\ \mathrm{bytes},\\
 \Delta B &= 25.00\%,\\
 h &= 64/256 = 25.00\%,\\
 B_{\mathrm{peak,resident}} &= 268{,}296{,}192\ \mathrm{bytes}
 < 256\ \mathrm{MiB}.
\end{align*}
All 757{,}334{,}016 logical returned bytes matched the declared demand.  The store
recorded 192 faults and 101 evictions.  Descriptive median route time was
1.06744~s and nearest-rank p95 was 2.13648~s, but timing is not promoted because
operating-system page-cache state was not controlled.  The report file SHA-256 is
\texttt{7046ec0d6c1f61a499cd4c1aea7f07440fac8479b3fbc8872f74167d88e8cce7};
its canonical digest is
\texttt{7bd35f13fca7ccbaeb1d264c547d53cd4817d86658b2a3026d96fca3fec4d3a3}.
Focused tests passed 3/3.  The bounded store was committed as
\texttt{bc088100}.

\subsection{Complete real Qwen route capture}

The generic debug frontend was first filtered to the internal
\texttt{ffn\_moe\_topk} tensors.  It confirmed that all 48 layer routers execute,
but its human-readable formatter replaced the middle two of every eight-expert
route with an ellipsis.  That diagnostic was preserved and invalidated for cache
use.  A small local callback was then compiled against the installed llama/ggml
interfaces.  It observes only the top-$k$ tensors, copies all eight \texttt{i32}
identifiers after evaluation, and rejects incomplete, duplicated or out-of-range
routes.

For the fixed minimal prompt \texttt{Hello}, the corrected diagnostic captured
all eight experts in all 48 layers.  It reported 48/48 complete layers, 384 valid
identifiers, eight distinct identifiers per layer, and no identifier outside
$[0,127]$.  The JSON route artifact SHA-256 is
\texttt{f9a0c6ee2ae269547dba62f4c9cc133c89988d52559ba0a82896d0a1492e66de}.
Its 242.26~s wall time is diagnostic and is not throughput evidence.  The native
capture mechanism is isolated in commit \texttt{03e3abee}.  This closes the gap
between the model's learned router and AION's exact expert byte addresses; the
next cache trace can use model-selected IDs rather than invented routes.

\subsection{Stride-correct prompt corpus and measured reuse}

Extending the callback to a 22-token business prompt exposed a further integrity
defect.  The first multi-token capture passed its range and uniqueness checks but
selected nearly every expert with almost no adjacent reuse.  A distribution audit
revealed that \texttt{ffn\_moe\_topk} is a strided eight-row view of a 128-row
argsort tensor.  Reading it as contiguous memory silently included non-top-$k$
rows.  That artifact was invalidated.  The corrected callback performs one
stride-aware eight-integer copy per token position.

The new corpus captured 22 routes in layers 0--46 and one route in layer 47.
The final-layer count is intentional: llama.cpp requests logits only for the last
prompt position and prunes the other final-layer positions.  Every route contained
eight unique IDs in $[0,127]$.  Across 987 adjacent route pairs, mean overlap was
3.1925 of eight experts and median overlap was three.  The histogram contained
61 zero-overlap pairs and 926 pairs with at least one reused expert.  The corpus
made 8{,}280 expert accesses but touched only 2{,}656 unique layer/expert
instances.  Corpus SHA-256:
\texttt{9468395059a85391c02080a631f92b999185b04202644a648df7d08cad082004}.

Replaying this genuine route in actual prompt layer-major order through the
256~MiB store produced:
\begin{align*}
 B_{\mathrm{logical,no\ cache}} &= 23{,}621{,}566{,}464\ \mathrm{bytes},\\
 B_{\mathrm{physical}} &= 7{,}576{,}731{,}648\ \mathrm{bytes},\\
 \Delta B &= 67.9245\%,\\
 h &= 5{,}624/8{,}280 = 67.9227\%,\\
 B_{\mathrm{peak,resident}} &= 268{,}406{,}784\ \mathrm{bytes}
 < 256\ \mathrm{MiB}.
\end{align*}
The store recorded 2{,}656 faults and 2{,}569 evictions.  Descriptive store time
was 61.64~s, with per-layer p50 1.70353~s and nearest-rank p95 2.17476~s.
These times are not promoted because page-cache state was uncontrolled and this
reader is not integrated with matrix computation.  In particular, the logical
no-cache denominator describes repeated addressed reads; it is not asserted to
equal stock llama.cpp's physical I/O.  The result proves substantial real-route
reuse and strict bounded residency, not end-to-end acceleration.  Report file
SHA-256:
\texttt{b3d186af4732e2adf91a20cb1424bbfb29782b7c27bc37084ecdb5889de9b6d6};
canonical digest:
\texttt{7b3e61ade808ca103d55dea77baa57fe8d6bbacc2e70f5302f42235074a0ffe9}.
Focused tests passed 3/3; implementation commit \texttt{acced8fd}.

\subsection{Direct packed CPU expert-compute boundary}

The next narrow gate tested whether the Mac CPU could compute directly on Qwen's
existing packed representation, avoiding FP16 expansion and Metal residency.
Layer~0 expert~21, selected by the real route, was read as 884{,}736 bytes of
Q4\_K gate weights, 884{,}736 bytes of Q4\_K up weights and 1{,}290{,}240 bytes
of Q6\_K down weights.  A native ggml graph executed the complete
$2048\rightarrow768\rightarrow2048$ SwiGLU expert.  Across 100 repetitions its
p50 was 0.058125~ms and p95 was 0.107084~ms, with finite output.  Artifact
SHA-256:
\texttt{7d1fd174e3652f203ce75bed639c406b86e1647894bb1737fd60633d0b03f263}.

The exact layer-0 route $[21,29,50,22,75,6,33,116]$ was then executed as eight
complete packed experts with summed outputs.  It consumed exactly 24{,}477{,}696
packed weight bytes.  Two independent 100-repetition runs produced the identical
checksum $-36.9675$:
\begin{table}[ht]
\centering
\begin{tabular}{lrrr}
\toprule
Run & p50 (ms) & p95 (ms) & Expert-only 48-layer p50 projection \\
\midrule
v1 & 0.427917 & 0.806666 & 48.6855 tokens/s \\
v2 & 0.416334 & 0.793791 & 50.0400 tokens/s \\
\bottomrule
\end{tabular}
\caption{Direct packed CPU timing for one real eight-expert Qwen layer route.}
\end{table}
Both runs passed the predeclared 0.694445~ms p50 boundary corresponding to an
expert-only 30-token/s projection.  V2 narrowly crossed the aspirational
0.416667~ms expert-only 50-token/s boundary; V1 did not.  The correct claim is a
replicated 48.69--50.04 token/s \emph{expert-compute ceiling at p50}, not
end-to-end model speed.  P95 implies only approximately 25.8--26.2 expert-only
tokens/s.  Attention, dense computation, router work, SD faults, transfers,
cache management and orchestration remain excluded.  V1 and V2 file hashes are
respectively
\texttt{32f9aeb40f579dcdc3b826cd0e9144764db35f6aa3be1962b3f535e278ea117d}
and
\texttt{e5fd56263fd0283c362aaef810d52752ad3b76a16d602266fd8eb9e3affaaf42}.
The native mechanism is isolated in commit \texttt{3617063a}.

\subsection{Representative-layer packed-compute audit}

The initial layer-0 timing was not assumed to generalise.  A follow-up
exploratory audit executed real token-0 routes at layers
$\{0,8,16,24,32,40,47\}$, with 200 repetitions per layer.  This audit also
caught and corrected an implementation assumption: Qwen's down projection is
not uniformly Q6\_K.  Exactly 24 layers use Q6\_K down weights and 24 use Q4\_K
down weights.  A first exploratory layer-16 attempt made with the incorrect
fixed-Q6 assumption produced a non-finite result and was rejected; the
manifest-driven tensor type restored finite results for every sampled layer.

The observed p50 range was 0.368709--0.430375~ms for eight experts in one
layer; p95 ranged from 0.682000 to 0.855583~ms.  The median sampled Q4\_K-down
layer cost was 0.369792~ms, while the median sampled Q6\_K-down cost was
0.423833~ms.  Applying those two descriptive medians to the actual 24/24 layer
split gives a 19.047~ms expert-only token projection, or 52.5017 tokens/s.
This crosses the aspirational 50-token/s boundary only as a stratified narrow
compute projection.

The result remains \texttt{EXPLORATORY\_NOT\_PROMOTED}.  The input activation
was synthetic; selected-expert gate coefficients were not applied; numerical
equivalence to llama.cpp was not tested; and routing, attention, dense work,
SD faults, cache control and orchestration were excluded.  It therefore does
not establish 52.5 end-to-end tokens/s.  It does strengthen the engineering
hypothesis that direct packed CPU expert arithmetic is not the limiting
physical barrier once the correct expert bytes are resident.  The SD-backed
report SHA-256 is
\texttt{f9c885d73c4f35741f7422ba1bb042fcd424ae545c398c1cc42acac13bdc3867}.

\subsection{First exact direct-packed reconstruction}

The correctness gate was then exercised on a real layer-0 token.  The stock
llama.cpp graph supplied the normalized 2{,}048-element layer input, selected
route $[62,21,75,87,36,34,126,47]$, eight normalized router coefficients and
the 2{,}048-element stock MoE output.  An independent graph read only those
eight packed expert slices from the SD-resident GGUF, evaluated the native
Q4\_K/Q6\_K matrix operations directly on the CPU, applied the captured router
coefficients and reduced the expert outputs.

All 2{,}048 reconstructed output values equalled the stock graph bit for bit:
maximum absolute error, mean absolute error and relative $L_2$ error were all
exactly zero.  This is the first proof that the direct-packed path is not merely
fast in a synthetic benchmark: for this real layer and token it reproduces the
reference MoE calculation exactly without expanding the expert weights to
FP16 or placing the complete expert bank in Metal memory.  The SD-backed report
SHA-256 is
\texttt{a9189f7d99dd0aa8cdad4b0057e9c10a38b62748f57abc44ecefe412756ab2bd}.

The acceptance boundary remains narrow.  One layer and one prompt token do not
establish all-layer, multi-token, step-logit or generated-token equivalence.
No end-to-end throughput claim is derived from this exactness result.  The next
gate must repeat the reconstruction over representative Q4\_K- and Q6\_K-down
layers and multiple real tokens before the mechanism is integrated into a
bounded decoder.

The representative-layer expansion then repeated the comparison at layers 16
and 32 on the same prompt, and at layer 8 on a distinct seven-word business
prompt.  This covered both Q4\_K and Q6\_K down projections.  Every one of the
8{,}192 compared floating-point values was bit-for-bit equal to the stock graph;
all four cases had zero maximum, mean and relative $L_2$ error.  The cohort
therefore passed the representative-layer exactness gate.  Its SD-backed report
SHA-256 is
\texttt{df3f57fb5419a621f96eb944dc597f5fa78e84cabf07cc287a25c823bd79d0a8}.
This broader result still does not substitute for every-layer, multi-token and
step-logit equivalence in an integrated decoder.

The gate was next expanded to the complete 48-layer stack in one stock capture
followed by independent reconstruction.  For every layer, AION read only the
selected eight packed expert slices, used the layer's actual Q4\_K or Q6\_K down
type, applied the captured normalized gates and compared the 2{,}048-value
result.  All 98{,}304 values were bit-for-bit equal: no layer had a nonzero
error, and aggregate maximum absolute, mean absolute and relative $L_2$ error
were zero.  The all-layer single-token exactness gate therefore passed.  The
SD-backed report SHA-256 is
\texttt{bbb85daf40b24f142ff1a89e5fac479914972bd2c5069e1f1d783deb4484fff4}.
This removes layer coverage as an immediate correctness uncertainty, but does
not yet prove multiple token positions, step logits, generated tokens or
end-to-end decoder speed.

\subsection{Next measured gates}

The next implementation stage is to combine the already bounded Qwen expert
reader with the now exact direct-packed calculation, without registering the
complete GGUF in Metal residency.  The immediate correctness gate is multi-token
reconstruction across token positions, followed by step-logit and generated-token
equivalence in the bounded decoder.  Only after those gates may persistent
buffers, decode/read/compute overlap and route-aware retention be judged by
end-to-end timing.  Exact-output candidates must preserve token
and step-logit equivalence; changed kernels or representations remain separately
quality-gated.  No claim of 50 tokens/s, production usability, or general larger-
model acceleration is made by the present results.
